\subsection{An algebraic two-challenge consistency checker}\label{sec:binary:checker}
We now combine the preceding lemmas into a small
\emph{algebraic checking theorem}. This is the
kernel-coordinate consistency component of the
proposed binary-field construction. The theorem
does not assert zero knowledge for the complete IOP
or its commitment scheme.

\begin{definition}[Affine consistency checks]\label{def:binary:checks}
Let the secret values $p,c,t$, $r,s$ have types
\[
 p,c,t\in L^{I_\ell},\qquad r,s\in L.
\]
After a challenge $\theta\in L^\ell$,
let $V_k=V_k(p,c,\theta)$ be as in
Definition~\ref{def:binary:affine-H}, and let
$U=(U_k)_k\in L^{I_\ell}$ be public.
After a second challenge $z\in L$, a public scalar
$A\in L$ is supplied. The \emph{algebraic
consistency checker} accepts precisely if
\begin{align}
 \Tr(U_k)&=0\quad\text{for all }k\in I_\ell,
 \label{eq:binary:check-traces}\\
 A+r+\sum_kW_k(z)t_k&=0,
 \label{eq:binary:check-first}\\
 A^2+s+\sum_kW_k(z)^2(U_k+V_k+t_k)&=0.
 \label{eq:binary:check-second}
\end{align}
The last two equations are affine in the
committed secret variables $p,c,t,r,s$ once
$\theta,z,U,A$ are fixed: $A^2$ and $W_k(z)^2$
are public scalars.
\end{definition}

\begin{definition}[Scalar and vector checking residuals]\label{def:binary:check-residuals}
Put
\[
 \delta:=s+r^2,\qquad
 e_k:=U_k+V_k+t_k^2+t_k
\]
and form the anchored residual polynomial
$F_{\delta,e}$ of Definition~\ref{def:binary:anchor-residual}.
\end{definition}

\begin{lemma}[First equation determines the public scalar]\label{lem:binary:first-determines}
For fixed $r,t,z$, equation~\eqref{eq:binary:check-first}
holds if and only if
\[
 A=r+\sum_kW_k(z)t_k.
\]
\end{lemma}
\begin{proof}
Move the other terms to the right side of
equation~\eqref{eq:binary:check-first}.
In characteristic two, additive inverses
equal the original values, so subtraction and
addition coincide. The resulting formula is
both necessary and sufficient.
\end{proof}

\begin{lemma}[The two affine checks imply an anchored residual root]\label{lem:binary:check-implies-root}
If equations~\eqref{eq:binary:check-first}
and~\eqref{eq:binary:check-second} both hold, then
\[
 \boxed{\delta+\sum_ke_kW_k(z)^2=0.}
\]
\end{lemma}
\begin{proof}
By Lemma~\ref{lem:binary:first-determines},
$A=r+\sum_kW_k(z)t_k$. Square this identity.
Characteristic-two additivity of squares
(Lemma~\ref{lem:binary:square}) gives
\[
 A^2=r^2+\sum_kW_k(z)^2t_k^2.
\]
Substitute it into equation~\eqref{eq:binary:check-second}:
\[
 0=r^2+s+
   \sum_kW_k(z)^2\bigl(t_k^2+U_k+V_k+t_k\bigr).
\]
The first two terms are $\delta=s+r^2$,
and the parenthesis is exactly $e_k$ by
Definition~\ref{def:binary:check-residuals}.
This is the claimed identity.
\end{proof}

\begin{lemma}[Honest trace check]\label{lem:binary:honest-traces}
Assume $c=c^p$ and choose
$U_k:=V_k+\wp(t_k)$.
Then the trace equations~\eqref{eq:binary:check-traces}
hold for every $\theta$ and every $t$.
\end{lemma}
\begin{proof}
By the definition of $c^p$ and
Definition~\ref{def:binary:transpose},
$D_c(\theta)=\Delta_p(\theta)$ for every $\theta$.
Corollary~\ref{cor:binary:masked-traces-consistent}
then yields $\Tr(U_k)=0$ for all $k$.
\end{proof}

\begin{lemma}[Honest affine checks]\label{lem:binary:honest-affine}
Assume $c=c^p$, $s=r^2$, and
$U_k=V_k+\wp(t_k)$ for every $k$.
For every $\theta,z$, choose
\[
 A:=r+\sum_kW_k(z)t_k.
\]
Then the two affine equations
\eqref{eq:binary:check-first} and
\eqref{eq:binary:check-second} hold.
\end{lemma}
\begin{proof}
The first equation holds by definition of $A$,
since $A+A=0$ in characteristic two.
The second equation follows from
$U_k+V_k+t_k=t_k^2$:
\[
 A^2+s+\sum_kW_k(z)^2(U_k+V_k+t_k)
 =A^2+r^2+\sum_kW_k(z)^2t_k^2.
\]
By the same squaring calculation used in
Lemma~\ref{lem:binary:check-implies-root},
$A^2=r^2+\sum_kW_k(z)^2t_k^2$.
The displayed sum therefore consists of two
identical copies of each term and equals zero.
\end{proof}

\begin{theorem}[Perfect algebraic completeness]\label{thm:binary:checker-complete}
If $c=c^p$, $s=r^2$, $U_k=V_k+\wp(t_k)$, and
$A=r+\sum_kW_k(z)t_k$, the checker accepts for
every choice of the challenges $(\theta,z)$.
\end{theorem}
\begin{proof}
Lemma~\ref{lem:binary:honest-traces} proves the
trace tests, and Lemma~\ref{lem:binary:honest-affine}
proves the two remaining affine equations.
Thus all tests in Definition~\ref{def:binary:checks}
hold.
\end{proof}

\begin{theorem}[Conditional anchored-check soundness]\label{thm:binary:checker-conditional}
Fix $\theta,p,c,t,r,s,U$ before a uniform
$z\overset{\$}{\leftarrow}L$ is sampled.
Suppose the trace checks hold and
$D_c(\theta)\neq\Delta_p(\theta)$.
Then, even if $A$ is chosen after seeing $z$,
the probability that both affine checks pass
is at most $n/q$.
\end{theorem}
\begin{proof}
Because the traces of all $U_k$ vanish while
$D_c(\theta)\neq\Delta_p(\theta)$,
Lemma~\ref{lem:binary:forced-residual}
shows that some $e_k$ is nonzero.
Thus $(\delta,e)\neq(0,0)$.
The residual values depend only on
$\theta,p,c,t,r,s,U$, which are fixed before $z$,
so the hypotheses of Theorem~\ref{thm:binary:anchor-root}
apply. If both affine checks pass,
Lemma~\ref{lem:binary:check-implies-root}
forces
$\delta+\sum_ke_kW_k(z)^2=0$.
The anchored-root bound gives probability
at most $n/q$, independently of how $A$
is selected after the challenge.
\end{proof}

\begin{theorem}[Two-challenge coefficient-consistency bound]\label{thm:binary:two-challenge}
Fix $p,c\in L^{I_\ell}$ with $c\neq c^p$ before
sampling uniform $\theta\in L^\ell$.
After seeing $\theta$, an arbitrary strategy
may select $t,r,s,U$ (and must satisfy all trace
checks), but these values must be fixed
\emph{before} an independent uniform
$z\overset{\$}{\leftarrow}L$ is sampled.
The strategy may choose $A$ after seeing $z$.
Then the probability that all checks in
Definition~\ref{def:binary:checks} accept is at most
\[
 \boxed{
 \min\left\{1,\frac{\ell+n}{q}\right\},
 \qquad q=2^n.
 }
\]
\end{theorem}
\begin{proof}
Let $E$ be the event
$D_c(\theta)=\Delta_p(\theta)$.
Since $c\neq c^p$ was fixed before $\theta$,
Corollary~\ref{cor:binary:mismatch-bound}
gives $\Pr[E]\leq\ell/q$.
On the complementary event $E^c$,
condition on any particular $\theta$ and on
the strategy's chosen values $t,r,s,U$
(and its private random tape), provided
the public trace checks pass.
These values are fixed before the fresh
challenge $z$, so
Theorem~\ref{thm:binary:checker-conditional}
bounds the conditional probability of
passing both affine checks by $n/q$.
If the trace checks fail, the acceptance
probability is zero. Averaging over these
conditional cases yields
\[
 \Pr[\text{accept}]
 \leq\Pr[E]+\Pr[\text{accept}\cap E^c]
 \leq\frac{\ell}{q}+\frac{n}{q}.
\]
Probabilities are also bounded by $1$.
This proves the stated minimum.
\end{proof}

\begin{lemma}[No even-order multiplicative subgroup in a binary field]\label{lem:binary:no-even-subgroup}
Every subgroup of $L^\times$ has odd order.
In particular, the odd-characteristic
pairing $x\leftrightarrow-x$ used in the
original KBFold word fold has no distinct
paired points in $L^\times$.
\end{lemma}
\begin{proof}
The group $L^\times$ has order
$|L|-1=2^n-1$, an odd integer.
The order of any subgroup divides
the order of the finite ambient group,
so every subgroup has odd order.
Furthermore $-x=x$ for every $x\in L$
because the characteristic is two.
Thus $x$ and $-x$ cannot form two distinct
points.
\end{proof}

\begin{remark}[Precise boundary of the result]
Theorems~\ref{thm:binary:checker-complete} and
\ref{thm:binary:two-challenge} concern only
the \emph{algebraic coefficient-consistency
checker}. They assume the stated timing
and do not incorporate a concrete
commitment-binding error, the remaining
ring-switching equations, the simulator
for the joint public messages $(U,A)$,
or any proof of security for the
odd-characteristic KBFold IOP.
The final lemma also proves that the
binary checker cannot simply reuse the
$\pm x$ multiplicative-domain fold.
These issues will be addressed in the
appropriate later chapters.
\end{remark}
